\section*{Research Methodology}
\addcontentsline{toc}{section}{Research Methodology}

This work adopts an interdisciplinary, systems-oriented qualitative methodology, synthesizing concepts from fields such as philosophy, systems theory, complexity science, history, law, and AI to construct an integrated conceptual framework. Rather than conducting empirical experiments, the research relies on conceptual analysis, comparative study, and logical reasoning to identify recurring structural patterns and systemic principles. The analytical process was highly iterative, organically refining ideas through continuous cycles of formulation, critique, and synthesis. Historical events, scientific developments, and cultural traditions are examined as case studies and comparative tools, rather than to assert historical equivalence. Advanced artificial intelligence systems supported the research and drafting process; a transparent, comprehensive disclosure of this human-AI co-development and authorship boundary is detailed in the preliminary pages. Emphasizing intellectual honesty and reproducibility, this framework distinguishes argument from evidence and is presented not as a definitive conclusion, but as an evolving matrix to stimulate future interdisciplinary inquiry.

\section*{AI Collaboration Disclosure}
This work was developed through a dynamic process of AI-human collaboration, where artificial intelligence served as an active collaborative partner in brainstorming, structural refinement, and conceptual synthesis. The core ideas, critical analysis, final editorial direction, and ultimate conceptual framework are entirely the author's own, and the author maintains full and sole responsibility for the final content of this paper.

